\section{Posterior inference}\label{sec:sampling}

Posterior inference is performed using a Markov Chain Monte Carlo (MCMC) algorithm. 
We extend the standard Metropolis-within-Gibbs sampler used in Bayesian additive regression trees to a reversible jump-within-Gibbs scheme \citep{ReversibleJump}.
This extension is required because the global--local horseshoe prior on the step heights breaks the conjugacy used in the original BART algorithm \citep{BARToriginal}. 
The original independent Gaussian prior allows analytical integration over the step heights. 
This enables separate updates of the tree structures and step heights.
We instead use a reversible jump Metropolis--Hastings step to jointly update both the tree structure and the step heights. 
Each MCMC iteration alternates between updating the prognostic and treatment effect forests using a Bayesian backfitting strategy \citep{Backfitting}.

In the remainder of this section:

\begin{enumerate}
    \item \textbf{Outer Gibbs sampler.}  
    We describe the Gibbs sampling strategy for updating the components of the decomposed AFT model. Separate forests are specified for the prognostic and treatment effect functions, and survival times are handled using data augmentation.

    \item \textbf{Reversible jump step.}  
    We detail the reversible jump Metropolis--Hastings step for updating a single tree. Our approach introduces a novel proposal mechanism for the step heights and local shrinkage parameters. These candidate values are drawn from pseudo-Gibbs updates informed by the parent node, which improves efficiency and mixing.

    \item \textbf{Full conditional updates.}  
    We describe the full conditional posterior updates for the step height parameters, conditional on a fixed tree structure. This includes efficient matrix-based sampling and closed-form updates for the global and local shrinkage parameters under the horseshoe prior. We highlight the connection to classical normal means models, which provides intuition for this structure.
\end{enumerate}

We provide details on computational cost, scalability, and empirical runtimes in the Supplementary Materials.

\subsection{Outer Gibbs sampler}\label{sec:outer_gibbs}
For both the prognostic regression function $f(x)$ and the treatment effect regression function $\tau(x)$, we specify separate Horseshoe Forests. 
We denote the ensemble for $f(x)$ by $\big(\mathcal{T}_j^f, \mathcal{H}_j^f\big)_{j=1}^{m_f}$ and for $\tau(x)$ by $\big(\mathcal{T}^\tau_j, \mathcal{H}^\tau_j\big)_{j=1}^{{m_\tau}}$, where $m_f$ and ${m_\tau}$ represent the number of trees in each respective forest. 
The outer Gibbs sampler is summarised in Algorithm~\ref{OuterGibbs}.

The trees in each forest are iteratively updated conditional on all other trees. 
This conditional updating step depends on the current model state only through the $n$ residuals, which are derived from the model in Eq.~\eqref{model}.
Specifically, the residuals for the $j$-th tree in the prognostic forest and treatment effect forest are given by:
\begin{subequations}\label{Residuals}
\begin{align}
\mathcal{R}_j^f &:= \log T^o - \sum_{J=1,\, J\neq j}^{m_f} g(X; \mathcal{T}_J^f, \mathcal{H}_J^f) - A \sum_{J=1}^{{m_\tau}} g(X; \mathcal{T}^\tau_J, \mathcal{H}^\tau_J), \label{Residuals_f} \\
\mathcal{R}^\tau_j &:= \log T^o - \sum_{J=1}^{m_f} g(X; \mathcal{T}_J^f, \mathcal{H}_J^f) - A \sum_{J=1,\, J\neq j}^{{m_\tau}} g(X; \mathcal{T}^\tau_J, \mathcal{H}^\tau_J), \label{Residuals_tau}
\end{align}
\end{subequations}
where $T^o$ denotes the observed or imputed survival times after data augmentation. 
We refer to $\mathcal{R}_j^f$ as the \textsl{prognostic residuals} and to $\mathcal{R}_j^\tau$ as the \textsl{treatment effect residuals}.


Algorithm~\ref{OuterGibbs} produces posterior draws of the conditional average treatment effect by evaluating differences of marginal conditional expectations under $A=1$ and $A=0$ given the current model parameters.
It does not generate paired potential outcomes $(T_i(1),T_i(0))$ for individual units.
All stochastic updates in the outer Gibbs sampler depend on the observed (or augmented) factual outcomes $T^o$ exclusively through the residuals defined in \eqref{Residuals}.
The procedure targets marginal conditional means and does not assume or attempt to recover a joint distribution for individual-level potential outcomes.
Inference on individual causal contrasts would require additional modelling assumptions on the joint distribution of the potential outcomes \citep{Oganisian2025Untangling, Li2023BayesianCausal}.





The censored survival times are augmented within the Gibbs sampler. 
At each iteration, these times are imputed by drawing from their conditional distribution. 
This conditional distribution is a truncated normal, with the lower bound given by the observed censoring time. 
The framework can be naturally extended to interval-censored data, in which case the truncated normal is bounded by the corresponding interval limits. 
The truncated normal distribution is centred at the current model-based prediction $\widehat{T}$ of the survival time, computed using \eqref{model}, and has variance equal to the current draw of $\sigma^2$.
Detailed derivations and further discussion can be found in \citep{TannerWong1987}.


The error variance $\sigma^2$ is sampled from its full conditional distribution.
The inverse-$\chi^2$ prior is conjugate.
This yields an inverse-gamma update based on the residuals $r_i = \log{T^o_i} - \log{\widehat{T}_i}$. %, where $\widehat{T}_i$ denotes the current prediction for individual $i$.
Detailed derivations can be found in \citet[Section 3.2 and 14.2]{Gelman2013BDA}.

We adopt the conventional binary treatment indicator $A$, coded as 1 for treated individuals and 0 for controls, to clearly communicate our model structure and causal assumptions. 
However, the proposed model can also accommodate the invariant parametrisation introduced by \citet{hahn2020bayesian}, which relies on a data-adaptive treatment coding scheme.  



\input AlgOuterGibbs


\subsection{Reversible jump step for updating a tree}\label{subsec:RJ_step}

The tree structure is jointly updated together with its corresponding set of step heights. 
We propose a new pair $(\mathcal{T}_j, \mathcal{H}_j)$, which is then accepted or rejected using a Metropolis--Hastings step with a reversible jump move \citep{ReversibleJump}.
The conditional posterior distribution of a single tree factorises as
\begin{equation}
    p\bigl((\mathcal{T}_j, \mathcal{H}_j) \mid \mathcal{R}_j, \sigma^2\bigr) \propto  p\bigl(\mathcal{H}_j \mid \mathcal{T}_j, \mathcal{R}_j, \sigma^2\bigr) \times p\bigl(\mathcal{T}_j \mid \mathcal{R}_j, \sigma^2\bigr).
\end{equation}
This factorisation enables a two-step proposal mechanism: first proposing a new tree structure $\mathcal{T}_j$, then proposing new step heights $\mathcal{H}_j$ conditional on the proposed structure. 
The overall proposal is evaluated using the reversible jump acceptance probability to ensure detailed balance.

\input tree_moves 

The proposal mechanism relies on three fundamental move types: \texttt{GROW}, \texttt{PRUNE}, and \texttt{CHANGE} \citep{BayesianCART}. 
Figure~\ref{NewTrees} schematically illustrates these moves and shows how the dimensionality of the step height parameter vector $\mathcal{H}_j$ changes accordingly. 
In a \texttt{GROW} move, the tree expands by adding a split, introducing additional step height parameters (and corresponding local horseshoe variance parameters). 
Conversely, a \texttt{PRUNE} move reduces the tree size by removing a split and eliminating corresponding parameters. 
A \texttt{CHANGE} move modifies an existing split without altering the number of terminal nodes.


Along with each move, new values for the step height parameters are proposed. 
The idea is to mimic a Gibbs sampler locally: propose new parameters informed by current values, as if the dimension were fixed. 
While this pseudo Gibbs proposal does not directly sample from the true posterior, it provides a reasonable approximation.
This is corrected using the Metropolis--Hastings acceptance step to leave the posterior invariant.

We illustrate the proposal mechanism for the \texttt{GROW} move.
Details for the \texttt{PRUNE} and \texttt{CHANGE} moves are provided in Supplementary Materials~\ref{appendix:comp_details} \citep{supplement}.
Suppose we perform a \texttt{GROW} move on leaf $L$.
Then two new child leaves $L$ and $L+1$ are grown.
We propose new values $h'_L$, $h'_{L+1}$, $\lambda'_L$, and $\lambda'_{L+1}$ by drawing from conditional distributions informed by the parent node parameters $h_L$ and $\lambda_L$.
These conditional distributions can be derived using Bayes' rule.
We do not need to propose a global shrinkage parameter $\tau$ since this dimensionality is not affected by the tree moves.


The horseshoe prior can be rewritten using an auxiliary variable representation \citep{HSGibbs}:
\begin{equation}\label{AuxilliaryHorseshoe}
\begin{alignedat}{2}
\tau^2 \mid \xi         &\sim \mathcal{IG}\left(\frac{1}{2}, \frac{1}{\xi}\right) &\quad \quad\text{and}\quad\quad  \lambda_\ell^2 \mid \nu_\ell &\sim \mathcal{IG}\left(\frac{1}{2}, \frac{1}{\nu_\ell}\right),\\
\xi   &\sim \mathcal{IG}\left(\frac{1}{2}, \frac{1}{\alpha^2}\right) &  \nu_\ell &\sim \mathcal{IG}\left(\frac{1}{2}, \frac{1}{\alpha^2}\right).
\end{alignedat}
\end{equation}
for $\ell = 1, \ldots, L$.
This representation allows all conditional distributions to be derived in closed form.
For notational simplicity, we write $\cdot$ in the conditioning statement to indicate dependence on all other model parameters and the data.
For the new node $L$, the new parameters are sequentially proposed as:
\begin{equation}
\begin{split}
\nu_L' \mid \cdot &\sim \mathcal{IG}\left(1, \frac{1}{\alpha^2} + \frac{1}{\lambda_L^2}\right),\\
\lambda_L'{}^2 \mid \cdot &\sim \mathcal{IG}\left(1, \frac{1}{\nu_L'} + \frac{h_L^2}{2\tau^2\omega}\right),\\
h_L' \mid \cdot &\sim \mathcal{N}\left( \frac{\bar{\mathcal{R}}_L'}{n_L' + \frac{1}{\tau^2\lambda_L'{}^2\omega}}, \frac{\sigma^2}{n_L' + \frac{1}{\tau^2\lambda_L'{}^2\omega}}\right),
\end{split}
\end{equation}
where $\bar{\mathcal{R}}_L'$ and $n_L'$ denote the sum and count of residuals mapped to child node $L$ by the proposed tree.  
The same procedure is applied to propose the parameters for node $L+1$.

A summary of the single-tree update is given in Algorithm~\ref{SingleTreeUpdate}.
Detailed acceptance ratios and derivations are provided in Supplementary Materials~\ref{appendix:comp_details}.
Related RJMCMC-based approaches for BART models in settings where step heights cannot be analytically integrated out have recently been studied by \citet{GeneralizedBART}.
While our modelling context and shrinkage structure differ, both approaches rely on joint proposals of tree structures and associated step heights within an RJMCMC framework.
\input AlgSingleTree





\subsection{Full conditional update of step height parameters}\label{subsection:full_conditional}

After performing a reversible jump Metropolis--Hastings step that acts locally on a specific terminal node of the tree $\mathcal{T}_j$, we carry out a full conditional update of the associated step height parameters $\mathcal{H}_j$. 
Conditional on the current tree structure, the dimensionality of the parameter space remains fixed, enabling efficient Gibbs sampling updates. 
Given a tree $\mathcal{T}$ and observed covariates $x_1, \ldots, x_n$, the $n$ observations are partitioned into $L$ step heights. 
This partition can be encoded by a design matrix $\mathcal{D} \in \mathbb{R}^{n \times L}$, where each row contains exactly one entry of 1 and zeros elsewhere, indicating the assigned step height. 
Specifically, if observation $i$ belongs to leaf $\ell$, then $\mathcal{D}_{i\ell} = 1$. 
For $\mathcal{D}=I_n$ this reduces to a normal means model.
By the valid partitions assumption, each row of $\mathcal{D}$ contains at least one non-zero entry, ensuring no empty leaves.

This structure allows us to reformulate the update as a regression problem in terms of the step heights $\smash{h = (h_1, \ldots, h_L)}$:
\begin{equation}
    \mathcal{R} \mid \mathcal{D}, h, \sigma^2 \sim \mathcal{N}\left(\mathcal{D}h, \sigma^2 I_n\right),
\end{equation}
where $\mathcal{R} \in \mathbb{R}^n$ denotes the residuals as defined in (\ref{Residuals}).

The conditional posterior of $h \in \mathbb{R}^L$ is given by \citep{lindley1972bayes}:
\begin{equation}
    h \sim \mathcal{N}_L\left(\Omega^{-1}\mathcal{D}^\top \mathcal{R}, \Omega^{-1}\right), \quad \text{where} \quad \Omega = \mathcal{D}^\top \mathcal{D} + \left[\omega\tau^2 \operatorname{diag}(\lambda_1^2, \ldots, \lambda_L^2)\right]^{-1}.
\end{equation}
Since $\mathcal{D}^\top \mathcal{D}$ is a diagonal matrix with leaf sizes on its diagonal, let $n_\ell$ denote the number of observations in leaf $\ell$ and $\bar{\mathcal{R}}_\ell$ the sum of residuals in leaf $\ell$. 
We can then express the conditional posterior of each $h_\ell$ as:
\begin{equation}
    h_\ell \mid \cdot \sim \mathcal{N}\left(\frac{\bar{\mathcal{R}}_\ell}{n_\ell + \frac{1}{\tau^2\lambda_\ell^2\omega}}, \frac{\sigma^2}{n_\ell + \frac{1}{\tau^2\lambda_\ell^2\omega}}\right).
\end{equation}

The conditional posteriors for the shrinkage parameters follow directly from the auxiliary parametrisation of the horseshoe prior, as presented in (\ref{AuxilliaryHorseshoe}):
\begin{equation}\label{PosteriorAuxiliary}
\begin{alignedat}{2}
\tau^2 \mid \cdot &\sim \mathcal{IG}\left(\frac{L+1}{2}, \frac{1}{\xi} + \frac{1}{2\omega} \sum_{\ell=1}^{L} \frac{h_\ell^2}{\lambda_\ell^2}\right)
&\quad\text{and}\quad&
\lambda_\ell^2 \mid \cdot \sim \mathcal{IG}\left(1, \frac{1}{\nu_\ell} + \frac{h_\ell^2}{2\tau^2\omega}\right), \\
\xi \mid \cdot &\sim \mathcal{IG}\left(1, \frac{1}{\alpha^2} + \frac{1}{\tau^2}\right)
&\quad\text{   }\quad&
\nu_\ell \mid \cdot \sim \mathcal{IG}\left(1, \frac{1}{\alpha^2} + \frac{1}{\lambda_\ell^2}\right).
\end{alignedat}
\end{equation}
All of these conditional distributions are inverse gamma, and together with the Gaussian update for $h$, enable efficient Gibbs sampling within each tree. 
This hierarchical structure facilitates scalable and adaptive updates while maintaining strong shrinkage properties, thus improving mixing and posterior exploration.












